\documentclass[10pt,a4paper]{article}
\usepackage[margin=0.56in]{geometry}
\usepackage{amsmath}
\usepackage{booktabs}
\usepackage{graphicx}
\usepackage{caption}
\captionsetup{hypcap=false}
\usepackage{microtype}
\usepackage{xcolor}
\usepackage{hyperref}
\usepackage{enumitem}
\usepackage{siunitx}
\setlist{nosep,leftmargin=*}
\setlength{\parindent}{0pt}
\setlength{\parskip}{2.5pt}
\setlength{\textfloatsep}{5pt}
\setlength{\floatsep}{5pt}
\setlength{\tabcolsep}{4.5pt}
\graphicspath{{../outputs/part_b/figures/}}
\hypersetup{colorlinks=true,urlcolor=blue!60!black,linkcolor=black}
\newcommand{\StudentName}{Student}
\InputIfFileExists{student-info.tex}{}{}

\title{SC4001 Programming Assignment 1\\Part B: Chest X-ray Classification}
\author{\StudentName}
\date{}

\begin{document}
\maketitle
\vspace{-1.6em}

\section{Experimental protocol}
The supplied partitions contain 4,708 training, 524 validation, and 624 test
images. Label 0 denotes normal/healthy and label 1 denotes pneumonia. Each
8-bit image was divided by 255 and then standardised using only the training
mean (0.571699) and population standard deviation (0.174098); no data
augmentation was used. The class counts were 1,214/3,494, 135/389, and 234/390
(normal/pneumonia) in the three partitions.

Every convolution used no bias and was followed by batch normalisation and
ReLU. Fully connected layers used bias. Following the course clarification,
dropout ($p=0.5$) was used only after the ReLUs of B1's hidden \texttt{fc,64}
and \texttt{fc,32} layers; B2 and B3 used no dropout. Each final \texttt{fc,1}
returned one unrestricted logit, passed directly to
\texttt{BCEWithLogitsLoss}. Adam, seed 42, a maximum of 20 epochs, patience 4,
and minimum improvement $10^{-4}$ were used throughout. Validation BCE loss
controlled early stopping and selected each candidate's best epoch checkpoint.
B1 learning rate and B2 batch size were then chosen by that checkpoint's lowest
validation classification error, with validation loss only as a deterministic
tie-break between hyperparameters. The test partition was evaluated once after
selection.

Training ran with deterministic PyTorch algorithms on one Tesla V100 32 GB
(PyTorch 2.5.1, CUDA 11.8). Sensitivity is
$\mathrm{TP}/(\mathrm{TP}+\mathrm{FN})$ for pneumonia and specificity is
$\mathrm{TN}/(\mathrm{TN}+\mathrm{FP})$ for normal cases.

\section{B1: AlexNet-like learning-rate search}
Batch size was fixed at 128. Table~\ref{tab:b1} reports the metrics of each
candidate's selected checkpoint.

\begin{center}
\small
\begin{tabular}{rrrrrr}
\toprule
Learning rate & Best epoch & Val. loss & Val. error & Sensitivity & Specificity \\
\midrule
$10^{-5}$ & 20 & 0.26925 & 0.09924 & 0.95887 & 0.73333 \\
$10^{-4}$ & 20 & 0.07148 & 0.02672 & 0.98458 & 0.94074 \\
$10^{-3}$ & 14 & 0.05353 & \textbf{0.01336} & 0.99486 & 0.96296 \\
$10^{-2}$ & 14 & 0.05979 & 0.01908 & 0.98972 & 0.95556 \\
$10^{-1}$ & 16 & 0.14681 & 0.04962 & 0.96144 & 0.91852 \\
\bottomrule
\end{tabular}
\captionof{table}{B1 validation results.}\label{tab:b1}
\end{center}

\begin{center}
\begin{minipage}[t]{0.49\textwidth}
\centering\includegraphics[width=\linewidth]{b1_validation_loss.png}
\end{minipage}\hfill
\begin{minipage}[t]{0.49\textwidth}
\centering\includegraphics[width=\linewidth]{b1_validation_accuracy.png}
\end{minipage}
\end{center}

The two smallest rates were still improving at epoch 20, especially
$10^{-5}$, so they converged too slowly under the common budget. The $10^{-1}$
curves were substantially worse, consistent with steps that were too large.
$10^{-3}$ attained the lowest validation error and was therefore selected;
$10^{-2}$ was close but slightly worse. On the test set, the selected model
obtained accuracy 0.82212, sensitivity 0.99744, and specificity 0.52991
(TP/TN/FP/FN = 389/124/110/1). Accuracy alone hides this asymmetric behaviour:
the model detected nearly every pneumonia image but incorrectly labelled 110
of 234 normal images as pneumonia. Sensitivity and specificity are therefore
required to distinguish the two error types, especially under class imbalance.

\section{B2: VGG-like batch-size search}
The B1 learning rate $10^{-3}$ was fixed. Table~\ref{tab:b2} reports the
selected checkpoint for each batch size.

\begin{center}
\small
\begin{tabular}{rrrrrr}
\toprule
Batch size & Best epoch & Val. loss & Val. error & Sensitivity & Specificity \\
\midrule
16  & 14 & 0.05321 & 0.02481 & 0.99486 & 0.91852 \\
32  & 8  & 0.06605 & 0.02481 & 0.97943 & 0.96296 \\
64  & 12 & 0.07744 & \textbf{0.02290} & 0.98458 & 0.95556 \\
128 & 6  & 0.12555 & 0.04389 & 0.97943 & 0.88889 \\
256 & 15 & 0.10644 & 0.03053 & 0.97172 & 0.96296 \\
\bottomrule
\end{tabular}
\captionof{table}{B2 validation results at learning rate $10^{-3}$.}\label{tab:b2}
\end{center}

\begin{center}
\begin{minipage}[t]{0.49\textwidth}
\centering\includegraphics[width=\linewidth]{b2_validation_loss.png}
\end{minipage}\hfill
\begin{minipage}[t]{0.49\textwidth}
\centering\includegraphics[width=\linewidth]{b2_validation_accuracy.png}
\end{minipage}
\end{center}

Batch size 64 had the lowest validation error at its loss-selected checkpoint
and was chosen. Batch size 16 reached a lower validation loss, but its
loss-selected checkpoint had a higher classification error; this illustrates
the distinction between epoch selection and batch-size selection. Sizes 16 and
32 were close, whereas 128 and 256 gave higher errors. The final VGG-like model
achieved test accuracy 0.84936, sensitivity 0.98462, and specificity 0.62393
(TP/TN/FP/FN = 384/146/88/6).

\section{B3: Compact ResNet-like model and comparison}
The ResNet-like model used learning rate $10^{-3}$ and the selected batch size
64. Its best checkpoint was epoch 14. It achieved validation accuracy 0.96374
and test accuracy 0.86699, with test sensitivity 0.98974 and specificity
0.66239 (TP/TN/FP/FN = 386/155/79/4).

\begin{center}
\small
\begin{tabular}{lrrrrrr}
\toprule
Architecture & Accuracy & Sensitivity & Specificity & Error & Parameters & FLOPs \\
\midrule
AlexNet-like & 0.82212 & 0.99744 & 0.52991 & 0.17788 & 99,569 & 39,063,616 \\
VGG-like     & 0.84936 & 0.98462 & 0.62393 & 0.15064 & 18,185 & 19,464,256 \\
ResNet-like  & \textbf{0.86699} & 0.98974 & \textbf{0.66239} & \textbf{0.13301} & 77,169 & 25,297,024 \\
\bottomrule
\end{tabular}
\captionof{table}{Test metrics and architecture costs. FLOPs count Conv2d and
Linear operations with two FLOPs per multiply-accumulate; batch normalisation,
ReLU, pooling, and residual additions are excluded.}\label{tab:b3}
\end{center}

The VGG-like network was smallest and cheapest, yet improved markedly over the
AlexNet-like model; repeated $3\times3$ convolutions provide a useful accuracy/
cost trade-off here. AlexNet-like had the most parameters and FLOPs but the
lowest accuracy and specificity, so greater size did not imply better
generalisation. Residual shortcuts can ease optimisation; here the deeper
ResNet-like network obtained the best test accuracy and specificity, exceeding
VGG-like by 1.76 percentage points in accuracy while using about 4.2 times as
many parameters and 1.3 times the stated FLOPs. All three models had lower test than
validation accuracy, driven mainly by false positives; this is consistent with
overfitting and/or a partition shift, but the experiment cannot distinguish
those causes. Results come from one seed, so small architecture differences
should not be treated as robust rankings.

\section{B4: Grad-CAM comparison}
Five pneumonia test images were sampled without replacement using seed 2026
(indices 13, 120, 241, 404, and 532). For each trained model, Grad-CAM used the
last activated convolutional feature map and gradients of the pneumonia logit;
the resulting non-negative map was bilinearly upsampled and independently
normalised to $[0,1]$ for overlay.

\begin{center}
\includegraphics[width=0.94\textwidth]{b4_gradcam_comparison.png}
\end{center}

AlexNet-like generally produced broad central/lower-lung activation. VGG-like
maps were fragmented and sometimes emphasised image edges. ResNet-like often
concentrated near the lower image boundary, with a central response in one
case. The architectures therefore did not identify a single consistent region
across all five cases. These
overlays are qualitative explanations of model sensitivity, not evidence that
the highlighted pixels are clinically valid pneumonia lesions. No lesion
annotations were available, per-image colour normalisation prevents magnitude
comparison, and edge/border attention may reflect acquisition artefacts.

\end{document}
